% !PRINTED 382
\weberpage{382}
%% [running head, not body text -- commented out 2026-08-07] Page furniture (running
%% head, folio, signature) is excluded from this edition by standing policy.
%% 382 \hfill Fifteenth Chapter. \hfill .§ 104. % [sic] stray period before section mark as printed

First of all, for every class $A$ whatever and for arbitrary
fundamental divisors $\delta, \delta_1, \delta_2$, we have:

\weq{7}{\chi(1, A) = 1,}
\weq{8}{\chi(\delta_1, A)\chi(\delta_2, A) = \chi(\delta_1\delta_2, A),}
\weq{9}{\chi(\delta, A_0) = 1,}

in which $\delta_1\delta_2$ is the symbolic product defined above, and $A_0$
denotes the principal class by which the number 1 is representable.

If for a class $A$ the characters $\chi(\delta_1, A)$, $\chi(\delta_2, A)$ $\dots$
$\chi(\delta_\lambda, A)$ are given, then by (4) and (8) \gesperrt{all} the $\chi(\delta, A)$ are determined.
Each of these characters, however, can be $= +1$ or $= -1$,
and so there result $2^\lambda$ possible determinations of the characters.
But these are not independent of one another.

For from (5) it follows that
\[
Dy^2 \equiv (2ax + by)^2 \ (\mathrm{mod}\ 4m),
\]
and from this it follows, if one assumes $m$ relatively prime to $a$, and consequently to $y$: % [sic] German source prints relaliv.
\[
(Dy^2, m) = (\Delta, m) = 1.
\]
From this it results that for every class $A$

\weq{10}{\chi(\Delta, A) = +1.}

Now to every fundamental divisor $\delta$ there can be determined a definite
\gesperrt{complementary} fundamental divisor $\delta'$, different from $\delta$, in such a way
that, in the sense of the symbolic multiplication,

\weq{11}{\delta\delta' = \Delta}

holds. The complement of $\delta'$ is again $\delta$, and every prime divisor
which divides $\delta$ and $\delta'$ at the same time must divide $Q$.

From this there follows, by (10), for every class $A$:

\weq{12}{\chi(\delta, A) = \chi(\delta', A),}

and accordingly the number of possible determinations of
the $\chi(\delta, A)$ is now only $2^{\lambda-1}$.

One unites into a \gesperrt{genus (Genus)} all classes of forms $A$ in which all the characters $\chi(\delta, A)$ agree,
and arrives at the result:

5. \gesperrt{The number of existing genera is
at most equal to $2^{\lambda-1}$.}

Whether the number of existing genera is really so
great is a deeper question, to which we shall return later.
